\clearpage
\section{Instruction Execution Model}\label{page:instruction-execution-model}

This section defines the common execution contract used by the instruction descriptions. An individual
instruction may add stricter rules, but it does not silently replace the rules below.

\subsection{Architectural Result Priority}
When one instruction could encounter more than one exceptional condition, the architecturally reported result is
the first applicable condition in the following order.

\begin{enumerate}
\item Instruction fetch and fetch translation.
\item Instruction framing and acquisition of the complete instruction record.
\item Opcode recognition.
\item Availability of every extension required by the decoded operation.
\item Fixed-field, reserved-field, and effective-address-form legality.
\item Privilege checks.
\item Selector, control-image, and other architectural-state validation.
\item Predicate evaluation for a conditional operation.
\item Every non-suppressed operand access in architectural access order.
\item Execution-stage integer and floating-point exceptions.
\end{enumerate}

A false predicate suppresses the conditional operation's target-memory and implicit stack accesses. It does not
suppress instruction framing, opcode and extension checks, encoding legality, privilege checks, or control-state
validation. Explicit operands are accessed in instruction operand order unless an instruction states another order.
For a memory-to-memory operation, the complete source read precedes access to the destination, so a source-access
failure has priority over a destination-access failure.

Implicit accesses use the order stated by the instruction. In particular, PUSH captures its source before accessing
the destination stack slot; POP reads the stack slot before validating or writing its destination; PUSHP uses listed
register order and POPP uses reverse listed order; CALL validates its target before accessing the return stack; RET
reads the return stack before validating its target; and the far call, far return, event return, and supervisor-entry
instructions use the step order in their instruction descriptions.

\subsection{Implicit Stack Operations}
Ordinary implicit stack operations capture the current SS and SP before their first stack access. Every stack slot is
64 bits and is addressed through the captured SS. The ranges below are expressed in terms of the captured
\texttt{old\_sp}; the applicable segment, translation, and complete-range checks finish before the operation commits.

\Needspace{3.4in}
\manualtablecaption{Ordinary Implicit Stack Operations}
\begin{manualtabular}{@{}>{\raggedright\arraybackslash}p{0.92in}>{\raggedright\arraybackslash}p{1.18in}>{\raggedright\arraybackslash}p{0.78in}>{\raggedright\arraybackslash}p{2.72in}@{}}
\toprule
\rowcolor{ManualHeaderFill}
\textbf{Operation} & \textbf{Stack range} & \textbf{Committed SP} & \textbf{Slot layout or value order}\\
\midrule
\texttt{PUSH}, \texttt{CALL} &
$[\texttt{old\_sp}-8,\texttt{old\_sp})$ &
\texttt{old\_sp}-8 &
one source value or the following-instruction return PC at \texttt{old\_sp}-8\\
\texttt{POP}, \texttt{RET} &
$[\texttt{old\_sp},\texttt{old\_sp}+8)$ &
\texttt{old\_sp}+8 &
one destination value or return PC read at \texttt{old\_sp}\\
\texttt{PUSHP}, \texttt{FPUSHP} &
$[\texttt{old\_sp}-16,\texttt{old\_sp})$ &
\texttt{old\_sp}-16 &
listed pair's second register at \texttt{old\_sp}-16 and first register at \texttt{old\_sp}-8\\
\texttt{POPP}, \texttt{FPOPP} &
$[\texttt{old\_sp},\texttt{old\_sp}+16)$ &
\texttt{old\_sp}+16 &
listed pair's second register from \texttt{old\_sp}, then first register from \texttt{old\_sp}+8\\
\texttt{LCALL} &
$[\texttt{old\_sp}-16,\texttt{old\_sp})$ &
\texttt{old\_sp}-16 &
return PC at \texttt{old\_sp}-16 and return CS image at \texttt{old\_sp}-8\\
\texttt{LRET} &
$[\texttt{old\_sp},\texttt{old\_sp}+16)$ &
\texttt{old\_sp}+16 &
return PC at \texttt{old\_sp} and return CS image at \texttt{old\_sp}+8\\
\bottomrule
\end{manualtabular}

The complete range and all source values are established before a push operation changes memory or SP. A pop
operation reads its complete range and completes any applicable destination validation before changing a destination
or SP. A successful operation commits all listed stack and register effects together; a preceding fault exposes none
of them.

Architectural event entry and ERET use the event-frame layout in the Architectural Event Processing Model.
SYSCALL and SYSRET keep their continuation in the user-return control-register bank described by the Privileged
Execution Model.

\subsection{Operand Evaluation and EA Defaults}
Operands are decoded in instruction order. Source reads complete before the final destination write unless an
atomic form says otherwise. Effective-address calculation may produce an address without reading memory.

Auto-update EA forms update a temporary operand-evaluation image. The architectural register update becomes
visible only when the instruction commits.

Ordinary forms accept at most one memory operand. The explicitly encoded memory-memory exceptions are
@MEMORY_MEMORY_EXCEPTIONS@. Repeat behavior is available only through the repeat instructions and body-eligibility
rules defined in Repeated Scalar Execution.

FLAGS and FFLAGS fields that an instruction does not mention remain unchanged. Operand evaluation follows
instruction operand order, and an instruction has one architectural commit point unless its description explicitly
defines partial completion.

\subsection{Shared Side-Effect Rules}
An instruction that faults before commit leaves destination registers, destination memory, and auto-update
register effects unchanged unless the instruction explicitly defines partial completion.

Atomic read-modify-write forms have a single architectural commit point for the memory update and any associated
register result. Cache, TLB, and system-state instructions describe their own completion and visibility rules.

Ordinary integer ALU forms use the common one-memory-operand rule. Compare and test forms may read two operands and
update FLAGS without storing their temporary result. EXTS and EXTZ additionally provide explicitly encoded
register-memory, memory-register, and memory-memory conversions; they preserve FLAGS unless a repeat rule observes a
derived value. Data movement, LEA, and segment-address operations also preserve FLAGS unless their descriptions say
otherwise.

Control transfers make PC and CS changes at a single control-transfer commit point. Atomic operations require an
addressable memory operand. An \texttt{AT=0} atomic operand must be naturally aligned and commits the
read-modify-write as one architectural update; an \texttt{AT=1} atomic operand raises \texttt{ADDRESS\_TYPE} and has
no byte-alignment check. TLB,
page-table context, and privileged control operations require supervisor privilege. Cache-maintenance operations
preserve FLAGS and define their own visibility contract. Approximate transcendental floating-point operations are
available only when CPUID reports \texttt{FPTRANSA} and marks the instruction's accuracy-contract ID present.
